% Self-contained report. Compile with pdfLaTeX in Overleaf (twice for contents).
\documentclass[12pt,a4paper]{article}
\usepackage[T1]{fontenc}
\usepackage[utf8]{inputenc}
\usepackage[english]{babel}
\usepackage[margin=2.5cm]{geometry}
\usepackage{xcolor}
\usepackage{listings}
\usepackage{booktabs}
\usepackage{microtype}
\usepackage[hidelinks]{hyperref}
\definecolor{listinggray}{RGB}{247,247,247}
\lstset{basicstyle=\ttfamily\footnotesize,backgroundcolor=\color{listinggray},
 frame=single,rulecolor=\color{gray},breaklines=true,breakatwhitespace=false,
 columns=fullflexible,keepspaces=true,showstringspaces=false,tabsize=4,
 captionpos=b,aboveskip=10pt,belowskip=10pt}
\setlength{\emergencystretch}{3em}
\begin{document}
\begin{titlepage}
\centering
\textbf{Ministry of Education, Culture and Research of the Republic of Moldova}\par\medskip
\textbf{Technical University of Moldova}\par\medskip
\textbf{Faculty of Computers, Informatics and Microelectronics}\par\medskip
\textbf{Department of Software Engineering and Automatic}\par
\vspace{2.5cm}
\textbf{ANDREI ELCISCEV}\par
\textbf{Group FAF-243}\par
\vspace{2cm}
{\Large\textbf{LABORATORY WORK No. 9}\par}
\medskip
{\Large\textbf{WATCHING SYSTEM CALLS WITH STRACE}\par}
\vspace{2cm}
\begin{flushright}
\textbf{Checked by:}\\
Rei\c{t}man Patricia
\end{flushright}
\vfill
\textbf{Chi\c{s}in\u{a}u, 2026}
\end{titlepage}
\pagenumbering{arabic}
\tableofcontents
\newpage
\section{Introduction}
This laboratory work investigates the interface between Linux applications and the
kernel using \texttt{strace}. The experiments trace a familiar command, compile a
small C program that writes and reads a file, and count selected system calls.
The traces connect file descriptors, byte counts, and return values with the
operations performed by each program.

All experiments were executed in the Linux Mint 22.3 VirtualBox VM, under the
\texttt{user} account, in \texttt{/home/user/os-lab9}. The directory was created
before the lab exercises. Observe answers are stored in \texttt{report.md} there.
The listings below contain actual captured commands and output in place of
screenshots. Long lines may wrap for readability. Timing values describe individual
measurements and are not performance benchmarks.

The supplied two-page assignment lists Part 3, ``Count one kind of call,'' in its
overview, but contains detailed instructions and Observe questions only for Parts
1 and 2. Section~\ref{sec:count} therefore records a supplementary experiment
covering that stated goal; no missing questions are assumed.

\section{Experimental Environment}

\begin{lstlisting}[caption={Working directory, kernel, compiler, and tracing tool}]
$ pwd; uname -sr; gcc --version | head -1; strace --version | head -1
/home/user/os-lab9
Linux 6.14.0-37-generic
gcc (Ubuntu 13.3.0-6ubuntu2~24.04) 13.3.0
strace -- version 6.8
\end{lstlisting}

The required compiler and tracing utility were already installed. The shell
environment was set to \texttt{LC\_ALL=C} for consistent English output.
Commands were executed through VirtualBox guest control, with standard output
captured to files rather than attached to an interactive terminal. This matters
for standard I/O buffering, discussed below.

\section{Part 1: Watching a Real Command}
\subsection{Tracing \texttt{ls -l}}
The option \texttt{-f} follows child processes if any are created. The expression
\texttt{-e trace=openat,read,write,close} limits the displayed calls to those four
names. Consequently, this trace does not show every call made by \texttt{ls}.
In particular, directory enumeration uses \texttt{getdents64}, which this filter
does not display. Opening libraries, reading configuration and user/group data,
and printing the listing also require kernel services.

\begin{lstlisting}[caption={Actual filtered trace of ls, including its normal output}]
$ strace -f -e trace=openat,read,write,close ls -l
openat(AT_FDCWD, "/etc/ld.so.cache", O_RDONLY|O_CLOEXEC) = 3
close(3)                                = 0
openat(AT_FDCWD, "/lib/x86_64-linux-gnu/libselinux.so.1", O_RDONLY|O_CLOEXEC) = 3
read(3, "\177ELF\2\1\1\0\0\0\0\0\0\0\0\0\3\0>\0\1\0\0\0\0\0\0\0\0\0\0\0"..., 832) = 832
close(3)                                = 0
openat(AT_FDCWD, "/lib/x86_64-linux-gnu/libc.so.6", O_RDONLY|O_CLOEXEC) = 3
read(3, "\177ELF\2\1\1\3\0\0\0\0\0\0\0\0\3\0>\0\1\0\0\0\220\243\2\0\0\0\0\0"..., 832) = 832
close(3)                                = 0
openat(AT_FDCWD, "/lib/x86_64-linux-gnu/libpcre2-8.so.0", O_RDONLY|O_CLOEXEC) = 3
read(3, "\177ELF\2\1\1\0\0\0\0\0\0\0\0\0\3\0>\0\1\0\0\0\0\0\0\0\0\0\0\0"..., 832) = 832
close(3)                                = 0
openat(AT_FDCWD, "/proc/filesystems", O_RDONLY|O_CLOEXEC) = 3
read(3, "nodev\tsysfs\nnodev\ttmpfs\nnodev\tbd"..., 1024) = 400
read(3, "", 1024)                       = 0
close(3)                                = 0
openat(AT_FDCWD, ".", O_RDONLY|O_NONBLOCK|O_CLOEXEC|O_DIRECTORY) = 3
close(4)                                = 0
close(4)                                = 0
openat(AT_FDCWD, "/etc/nsswitch.conf", O_RDONLY|O_CLOEXEC) = 4
read(4, "# /etc/nsswitch.conf\n#\n# Example"..., 4096) = 569
read(4, "", 4096)                       = 0
close(4)                                = 0
openat(AT_FDCWD, "/etc/passwd", O_RDONLY|O_CLOEXEC) = 4
read(4, "root:x:0:0:root:/root:/bin/bash\n"..., 4096) = 2876
close(4)                                = 0
close(4)                                = 0
close(4)                                = 0
openat(AT_FDCWD, "/etc/group", O_RDONLY|O_CLOEXEC) = 4
read(4, "root:x:0:\ndaemon:x:1:\nbin:x:2:\ns"..., 4096) = 1148
close(4)                                = 0
close(3)                                = 0
openat(AT_FDCWD, "/etc/localtime", O_RDONLY|O_CLOEXEC) = 3
read(3, "TZif2\0\0\0\0\0\0\0\0\0\0\0\0\0\0\0\0\0\0\0\0\0\0\20\0\0\0\0"..., 4096) = 2390
read(3, "TZif2\0\0\0\0\0\0\0\0\0\0\0\0\0\0\0\0\0\0\0\0\0\0\20\0\0\0\0"..., 4096) = 1491
close(3)                                = 0
write(1, "total 8\ndrwxrwxr-x 2 user user 4"..., 113total 8
drwxrwxr-x 2 user user 4096 Nov  2 16:17 evidence
-rw------- 1 user user 1742 Nov  2 16:17 experiment.sh
) = 113
close(1)                                = 0
close(2)                                = 0
+++ exited with 0 +++
\end{lstlisting}

The apparent mixing of the directory listing into a \texttt{write} line occurs
because normal output and \texttt{strace}'s diagnostic stream were captured together.
It is a recording effect, not an invalid system call.

\subsection{Observe: Printing to the Screen}
\textbf{Question:} Which system call does \texttt{ls} use to print to the screen?
To which descriptor does it write?

\textbf{Answer:} It uses \texttt{write}, writing to descriptor \texttt{1}, standard
output. In the captured run, the trace shows a 113-byte write to descriptor 1.
When run interactively, standard output normally points to the terminal; in this
experiment it was redirected into the evidence file.

\subsection{Rough Number of Calls}

\begin{lstlisting}[caption={Summary restricted to the four selected system calls}]
$ strace -f -c -e trace=openat,read,write,close ls -l
total 32
drwxrwxr-x 2 user user  4096 Nov  2 16:17 evidence
-rw------- 1 user user  1742 Nov  2 16:17 experiment.sh
-rw-r--r-- 1 user user    14 Nov  2 16:17 note.txt
-rwxrwxr-x 1 user user 16272 Nov  2 16:17 syscalls
-rw-rw-r-- 1 user user   663 Nov  2 16:17 syscalls.c
% time     seconds  usecs/call     calls    errors syscall
------ ----------- ----------- --------- --------- ----------------
 52.71    0.001564          97        16           close
 23.19    0.000688          62        11           read
 21.30    0.000632          63        10           openat
  2.80    0.000083          83         1           write
------ ----------- ----------- --------- --------- ----------------
100.00    0.002967          78        38           total
\end{lstlisting}

This measurement reports 38 selected calls: 10 \texttt{openat}, 11 \texttt{read},
16 \texttt{close}, and one \texttt{write}. This is a filtered count, not the total
number of system calls. A second measurement without a syscall filter gives the
full summary:

\begin{lstlisting}[caption={Unfiltered summary of system calls made by ls}]
$ strace -f -c ls -l
total 32
drwxrwxr-x 2 user user  4096 Nov  2 16:17 evidence
-rw------- 1 user user  1742 Nov  2 16:17 experiment.sh
-rw-r--r-- 1 user user    14 Nov  2 16:17 note.txt
-rwxrwxr-x 1 user user 16272 Nov  2 16:17 syscalls
-rw-rw-r-- 1 user user   663 Nov  2 16:17 syscalls.c
% time     seconds  usecs/call     calls    errors syscall
------ ----------- ----------- --------- --------- ----------------
 15.00    0.002673         243        11           read
 14.68    0.002617         163        16           close
 14.47    0.002580         151        17           mmap
 10.27    0.001831         183        10           openat
  7.35    0.001310         262         5           statx
  5.69    0.001015         338         3           newfstatat
  5.22    0.000931          71        13           fstat
  3.78    0.000674         674         1           execve
  3.32    0.000592         592         1           getrandom
  3.10    0.000553         110         5         5 lgetxattr
  3.05    0.000544         136         4         4 connect
  2.05    0.000365          73         5           mprotect
  2.01    0.000359         119         3           brk
  1.74    0.000311         311         1           arch_prctl
  1.51    0.000269          53         5           listxattr
  1.46    0.000260          65         4           socket
  1.35    0.000240         120         2         2 statfs
  0.74    0.000132         132         1           write
  0.72    0.000129         129         1           munmap
  0.51    0.000091          30         3           lseek
  0.47    0.000084          42         2         2 access
  0.39    0.000070          35         2           pread64
  0.30    0.000054          54         1           rseq
  0.30    0.000053          53         1           set_tid_address
  0.29    0.000051          51         1         1 ioctl
  0.12    0.000022          11         2           getdents64
  0.06    0.000011          11         1           prlimit64
  0.02    0.000003           3         1           set_robust_list
------ ----------- ----------- --------- --------- ----------------
100.00    0.017824         146       122        14 total
\end{lstlisting}

The unfiltered measurement reports 122 calls, including two \texttt{getdents64}
calls. Thus, a simple long listing makes roughly a hundred calls in this setup.
The directory had more files by the time of these summaries than during the first
trace because the program had been created and compiled. Counts depend on directory
contents, environment, libraries, and output destination. The 14 errors in the
summary are individual failed operations, including optional resource probes;
they do not mean that \texttt{ls} itself failed.

\section{Part 2: A Program Using File Descriptors}
\subsection{Source Code}
The supplied program was saved as \texttt{syscalls.c}. It opens \texttt{note.txt},
writes a message, closes the descriptor, opens the file again, and reads the message
back. File contents are transferred with \texttt{write} and \texttt{read};
\texttt{printf} reports the results on standard output.

\begin{lstlisting}[caption={Source code compiled in the VM}]
$ cat syscalls.c
/* syscalls.c - open, write, read and close. Lab 9 supplied program. */
#include <fcntl.h>
#include <unistd.h>
#include <string.h>
#include <stdio.h>

int main(void) {
    const char *msg = "hello, kernel\n";
    int fd = open("note.txt", O_CREAT | O_WRONLY | O_TRUNC, 0644);
    if (fd < 0) { perror("open"); return 1; }
    printf("open returned descriptor %d\n", fd);
    ssize_t n = write(fd, msg, strlen(msg));
    printf("write sent %zd bytes\n", n);
    close(fd);
    fd = open("note.txt", O_RDONLY);
    char buf[64];
    n = read(fd, buf, sizeof(buf) - 1);
    buf[n] = '\0';
    printf("read got %zd bytes: %s", n, buf);
    close(fd);
    return 0;
}
\end{lstlisting}

\subsection{Compilation and Execution}

\begin{lstlisting}[caption={Compilation, execution, and inspection of the written file}]
$ gcc -Wall -Wextra -o syscalls syscalls.c; ./syscalls; wc -c note.txt; cat note.txt
open returned descriptor 3
write sent 14 bytes
read got 14 bytes: hello, kernel
14 note.txt
hello, kernel
\end{lstlisting}

Compilation with \texttt{-Wall -Wextra} produced no diagnostics. The message
contains 13 visible characters plus a newline, so its length is 14 bytes.
\texttt{strlen} excludes the terminating null byte. The read requests up to 63
bytes and returns 14, the number actually available in the file.

\subsection{Inspecting the Actual Calls}
The required filtered trace was executed directly. For a more readable recording,
it was repeated with \texttt{-o}, which places the trace in a separate file while
leaving the program's standard output separate. Both outputs are reproduced below.
The unfiltered command \texttt{strace ./syscalls} was also run and saved in
\texttt{evidence/10\_full\_trace.txt} in the VM.

\begin{lstlisting}[caption={Program output followed by the separately captured filtered trace}]
$ strace -o evidence/syscalls-clean.trace -e trace=openat,write,read,close ./syscalls
open returned descriptor 3
write sent 14 bytes
read got 14 bytes: hello, kernel
$ cat evidence/syscalls-clean.trace
openat(AT_FDCWD, "/etc/ld.so.cache", O_RDONLY|O_CLOEXEC) = 3
close(3)                                = 0
openat(AT_FDCWD, "/lib/x86_64-linux-gnu/libc.so.6", O_RDONLY|O_CLOEXEC) = 3
read(3, "\177ELF\2\1\1\3\0\0\0\0\0\0\0\0\3\0>\0\1\0\0\0\220\243\2\0\0\0\0\0"..., 832) = 832
close(3)                                = 0
openat(AT_FDCWD, "note.txt", O_WRONLY|O_CREAT|O_TRUNC, 0644) = 3
write(3, "hello, kernel\n", 14)         = 14
close(3)                                = 0
openat(AT_FDCWD, "note.txt", O_RDONLY)  = 3
read(3, "hello, kernel\n", 63)          = 14
close(3)                                = 0
write(1, "open returned descriptor 3\nwrite"..., 80) = 80
+++ exited with 0 +++
\end{lstlisting}

The initial opens and read belong to dynamic-loader setup. The application then
opens \texttt{note.txt} with creation, write-only, and truncation flags and obtains
descriptor 3. Its \texttt{write(3,...)} returns 14. Closing the file frees the
descriptor, so opening it again can return 3. The subsequent \texttt{read} returns
the same 14 bytes, followed by another successful close.

Although the C source calls \texttt{open}, the C library implements that operation
through the Linux \texttt{openat} syscall here. \texttt{AT\_FDCWD} resolves a
relative pathname against the process's current working directory.

\subsection{Observe: Why Descriptor 3?}
\textbf{Question:} The descriptor returned by the first \texttt{openat} is 3, not
0, 1, or 2. Why are those three already taken?

\textbf{Answer:} The process inherits three standard descriptors: 0 is standard
input, 1 is standard output, and 2 is standard error. At the point where
\texttt{note.txt} is opened, these are occupied and 3 is the lowest available
descriptor. The temporary loader descriptor 3 has already been closed.

\subsection{Checking the Result}

\begin{lstlisting}[caption={Checks of the actual byte count and exact file contents}]
$ test "$(wc -c < note.txt)" -eq 14 && printf "PASS: note.txt contains 14 bytes\n"; printf "hello, kernel\n" | cmp - note.txt && printf "PASS: file content matches the written message\n"
PASS: note.txt contains 14 bytes
PASS: file content matches the written message
\end{lstlisting}

The original teaching program was retained. It checks the first open but does not
check the second open, the read result, or partial writes. In production code,
these must be checked before using \texttt{buf[n]} or assuming the whole message
was written. Both operations succeeded in the recorded experiment.

\section{Part 3: Counting One Kind of Call}
\label{sec:count}
This supplementary experiment uses \texttt{-c} to summarize calls and
\texttt{-e trace=write} to count only writes.

\begin{lstlisting}[caption={Counting write calls made by the program}]
$ strace -c -e trace=write ./syscalls
open returned descriptor 3
write sent 14 bytes
read got 14 bytes: hello, kernel
% time     seconds  usecs/call     calls    errors syscall
------ ----------- ----------- --------- --------- ----------------
100.00    0.000061          30         2           write
------ ----------- ----------- --------- --------- ----------------
100.00    0.000061          30         2           total
\end{lstlisting}

There are two \texttt{write} calls: one sends 14 bytes to \texttt{note.txt} through
descriptor 3, and one sends 80 bytes of status messages to standard output through
descriptor 1. The three \texttt{printf} calls do not necessarily produce three
kernel calls. Because output is captured through a non-terminal stream, the C
library buffers those messages and flushes them together at program exit.
Interactive terminal execution can produce a different count because standard
output is normally line buffered there.

\begin{lstlisting}[caption={Summary of the four selected calls for the program}]
$ strace -c -e trace=openat,write,read,close ./syscalls
open returned descriptor 3
write sent 14 bytes
read got 14 bytes: hello, kernel
% time     seconds  usecs/call     calls    errors syscall
------ ----------- ----------- --------- --------- ----------------
 60.28    0.000900         225         4           openat
 18.82    0.000281         140         2           write
 16.14    0.000241          60         4           close
  4.76    0.000071          35         2           read
------ ----------- ----------- --------- --------- ----------------
100.00    0.001493         124        12           total
\end{lstlisting}

The selected total is 12 calls: four \texttt{openat}, two \texttt{read}, two
\texttt{write}, and four \texttt{close}. These include loader activity as well as
the program's own file operations. The six calls specifically operating on
\texttt{note.txt} are two opens, one write, one read, and two closes.

\section{Results and Discussion}
The traces demonstrate that visible program actions can involve more kernel calls
than the source suggests. A directory listing loads libraries, obtains metadata,
looks up user and group names, and enumerates directory entries before printing.
Filtering helps isolate relevant calls but must not be mistaken for a complete
count.

The C program confirmed the relationship between descriptors and open files.
The same descriptor number was reused after closing the file; successful writes
and reads returned the number of transferred bytes. The correspondence between
\texttt{open} and \texttt{openat}, and between multiple \texttt{printf} calls and
one buffered \texttt{write}, also showed why library calls and system calls are
not always in one-to-one correspondence.

\section{Conclusion}
This laboratory work provided practical experience tracing Linux system calls,
interpreting descriptors and return values, and counting selected operations.
The program successfully wrote and read back the 14-byte message. The observations
explain standard descriptors, descriptor reuse, library startup activity, and the
effect of output buffering on measured calls. Actual command transcripts preserve
the experimental evidence without external screenshots or other report assets.
\end{document}
